\chapter*{Demonstration Mode}
\section*{Conditions}
This file was generated to see the results of three different algorithms solving a random generated 0/1 Knapsack Problem.
The problem has specific conditions everytime the \verb|knapsack| program is run.

\begin{itemize}
    \item The knapsack has a maximum size of 15.
    \item The are only 6 items.
    \item Each value for the items ranges between 0 to 20.
    \item Each cost  for the items ranges between 0 to 7.
\end{itemize}

\chapter*{Algorithms}

\section*{Dynamic Programming Algorithm}

This algorithm uses a table to get to the optimal solution for the generated 0/1 Knapsack Problem. The strategy here is to solve
subproblems, write down the optimal solution for that subproblem on a table and by using those solutions build the solution 
for the original problem. It always find the optimal solution for a 0/1 Knapsack Problem, as explained in course of Operations Research.

\section*{Basic Greedy Algorithm}

This algorithm focuses on trying to put inside the knapsack the most valuable item, if it does not fit in, tries with the next
most valueable item until the are no more items available.

\section*{Proportional Greedy Algorithm}

It has the same strategy as the Basic Greedy, but it focuses on trying to put inside items based by their proportion between its
value and cost. The item with highest proportion goes first, if it does not fit, goes to the next item with the highest proportion.